\chapter{Chapter09 — Noise Operator Exercises}
%%% generated by the blueprint pipeline from FABL.Chapter09.NoiseOperatorExercises
\section{Overview}

Declarations of the Lean module \texttt{FABL.Chapter09.NoiseOperatorExercises}.

\section{Declarations}

\begin{theorem}[Expectation on a sign cube by coordinates]
\lean{FABL.expect_snoc}
\label{FABL.expect_snoc}
\leanok
\uses{FABL.Sign, FABL.finSnocEquiv}
\difficulty{3}
Let $n$ be a nonnegative integer and let $f:\{-1,1\}^{n+1}\to\mathbb{R}$. The uniform
mean of $f$ on the $(n+1)$-dimensional sign cube equals the mean, over uniform
$x\in\{-1,1\}^n$ and an independent uniform sign $a\in\{-1,1\}$, of $f$ evaluated at the
vector obtained by appending $a$ to $x$.
\end{theorem}

\begin{theorem}[Product expectation by conditioning on the last coordinate]
\lean{FABL.pmfExpectation_independentProductPMF_succ}
\label{FABL.pmfExpectation_independentProductPMF_succ}
\leanok
\uses{FABL.finSnocEquiv, FABL.finSnocEquiv_apply, FABL.independentProductPMF, FABL.independentProductPMF_apply, FABL.pmfExpectation}
\difficulty{4}
Let $\Omega$ be a finite set, let $n$ be a natural number, let $p_0,\ldots,p_n$ be
probability mass functions on $\Omega$, and let $f:\Omega^{n+1}\to\mathbb{R}$. The
expectation of $f$ under the independent product distribution equals the expectation
over the last coordinate of the conditional expectation over the first $n$ coordinates:
$\mathbb{E}[f]=\mathbb{E}_{a\sim
p_n}\bigl[\mathbb{E}_{x\sim\bigotimes_{j=0}^{n-1}p_j}f(x_0,\ldots,x_{n-1},a)\bigr]$.
\end{theorem}

\begin{definition}[Even part in the last coordinate]
\lean{FABL.lastEvenPart}
\label{FABL.lastEvenPart}
\leanok
For a function $f:\{-1,1\}^{n+1}\to\mathbb R$, define its even part in the last
coordinate to be the function on $\{-1,1\}^n$ given by $x\mapsto
\bigl(f(x,1)+f(x,-1)\bigr)/2$.
\end{definition}

\begin{definition}[Odd part in the last coordinate]
\lean{FABL.lastOddPart}
\label{FABL.lastOddPart}
\leanok
For a function $f:\{-1,1\}^{n+1}\to\mathbb R$, define its odd part in the last
coordinate to be the function on $\{-1,1\}^n$ given by $x\mapsto
\bigl(f(x,1)-f(x,-1)\bigr)/2$.
\end{definition}

\begin{theorem}[Reconstruction from the last coordinate]
\lean{FABL.lastCoordinate_reconstruct}
\label{FABL.lastCoordinate_reconstruct}
\leanok
\uses{FABL.Sign, FABL.lastEvenPart, FABL.lastOddPart, FABL.signValue}
\difficulty{4}
Let $n\ge 0$, let $f:\{−1,1\}^{n+1}\to\mathbb R$, let $x\in\{−1,1\}^n$, and let
$a\in\{−1,1\}$. Then $f(x,a)=f_{\mathrm{even}}(x)+a f_{\mathrm{odd}}(x)$, where
$f_{\mathrm{even}}$ and $f_{\mathrm{odd}}$ are the averages and half-differences of the
two values obtained by setting the last coordinate to $1$ and $-1$, respectively.
\end{theorem}

\begin{theorem}[Noise operator on an appended coordinate]
\lean{FABL.noiseOperator_snoc}
\label{FABL.noiseOperator_snoc}
\leanok
\uses{FABL.Sign, FABL.lastCoordinate_reconstruct, FABL.lastEvenPart, FABL.lastOddPart, FABL.noiseKernel, FABL.noiseOperator, FABL.noiseOperator_apply_eq_pmfExpectation, FABL.pmfExpectation, FABL.pmfExpectation_add, FABL.pmfExpectation_const, FABL.pmfExpectation_const_mul, FABL.pmfExpectation_coordinateNoisePMF_signValue, FABL.pmfExpectation_independentProductPMF_succ, FABL.pmfExpectation_mul_const, FABL.signValue}
\difficulty{5}
Let $n$ be a natural number, let $-1\le\rho\le1$, let $f$ be a real-valued function on
the $(n+1)$-dimensional sign cube, let $x$ be a point of the $n$-dimensional sign cube,
and let $a\in\{ -1,1\}$. The noise operator with parameter $\rho$, applied to $f$ at the
point obtained by appending $a$ to $x$, equals the noise operator applied to the even
part of $f$ in its last coordinate at $x$, plus $\rho a$ times the noise operator
applied to the odd part of $f$ in its last coordinate at $x$.
\end{theorem}

\begin{definition}[Last-coordinate fiber norm]
\lean{FABL.lastFiberLpNorm}
\label{FABL.lastFiberLpNorm}
\leanok
\uses{FABL.Sign, FABL.uniformLpNorm}
For a real number $q$ and a function $f:\{-1,1\}^{n+1}\to\mathbb R$, define a function
on $\{-1,1\}^n$ by assigning to each $x$ the uniform $L^q$ norm of the two-value
function $a\mapsto f(x,a)$ on $a\in\{-1,1\}$.
\end{definition}

\begin{theorem}[Uniform norm over the last coordinate]
\lean{FABL.uniformLpNorm_eq_lastFiberLpNorm}
\label{FABL.uniformLpNorm_eq_lastFiberLpNorm}
\leanok
\uses{FABL.Sign, FABL.expect_snoc, FABL.lastFiberLpNorm, FABL.uniformLpNorm, FABL.uniformLpNorm_nonneg, FABL.uniformLpNorm_rpow}
\difficulty{4}
For a natural number $n$, a real number $q>0$, and a function
$f:\{-1,1\}^{n+1}\to\mathbb{R}$, define $g:\{-1,1\}^{n}\to\mathbb{R}$ by taking, at each
$x$, the uniform $L^q$ quantity of the two values $f(x,-1)$ and $f(x,1)$. Then the
uniform $L^q$ quantity of $f$ on $\{-1,1\}^{n+1}$ equals the uniform $L^q$ quantity of
$g$ on $\{-1,1\}^{n}$.
\end{theorem}

\begin{theorem}[Monotonicity of the uniform norm]
\lean{FABL.uniformLpNorm_le_of_abs_le}
\label{FABL.uniformLpNorm_le_of_abs_le}
\leanok
\uses{FABL.uniformLpNorm}
\difficulty{3}
Let $\Omega$ be a finite set, let $q>0$, and let $f,g:\Omega\to\mathbb R$. If
$|f(x)|\le |g(x)|$ for every $x\in\Omega$, then the uniform $L^q$ norm of $f$ is at
most that of $g$.
\end{theorem}

\begin{theorem}[Squared norm of the last-coordinate parts]
\lean{FABL.uniformLpNorm_two_lastParts_sq}
\label{FABL.uniformLpNorm_two_lastParts_sq}
\leanok
\uses{FABL.Sign, FABL.expect_snoc, FABL.lastCoordinate_reconstruct, FABL.lastEvenPart, FABL.lastOddPart, FABL.signValue, FABL.uniformLpNorm, FABL.uniformLpNorm_two_sq_eq_expect_sq}
\difficulty{4}
For every real-valued function $f$ on the $(n+1)$-dimensional sign cube, let $f_+$ and
$f_-$ be its restrictions to the last coordinate equal to $1$ and $-1$, respectively,
and define $f_{\mathrm{even}}=(f_++f_-)/2$ and $f_{\mathrm{odd}}=(f_+-f_-)/2$ on the
$n$-dimensional sign cube. Then
$\|f\|_2^2=\|f_{\mathrm{even}}\|_2^2+\|f_{\mathrm{odd}}\|_2^2$, where each norm is
uniform on its respective cube.
\end{theorem}

\begin{theorem}[Norm of a square root]
\lean{FABL.uniformLpNorm_sqrt_sq}
\label{FABL.uniformLpNorm_sqrt_sq}
\leanok
\uses{FABL.uniformLpNorm, FABL.uniformLpNorm_abs_rpow}
\difficulty{3}
Let $\Omega$ be a finite set, let $q>0$, and let $H:\Omega\to\mathbb{R}$ be nonnegative.
Then the square of the uniform $q$-norm of the function $x\mapsto\sqrt{H(x)}$ equals the
uniform $(q/2)$-norm of $H$.
\end{theorem}

\begin{theorem}[One bit two q sharp of two le]
\lean{FABL.oneBit_twoQ_sharp_of_two_le}
\statementsource{Analysis-of-Boolean-Functions-by-Ryan-ODonnell}{
  pdf-ffbe9c7b3306-p257-b003,
  pdf-ffbe9c7b3306-p257-b004
}
\label{FABL.oneBit_twoQ_sharp_of_two_le}
\leanok
\uses{FABL.Sign, FABL.oneBit_twoQ_sharp, FABL.signValue, FABL.uniformLpNorm, FABL.uniformLpNorm_two_affineSign}
\difficulty{4}
% dataset-pipeline: stub -- prose pending build_informal_statements.py
Statement pending regeneration.
\end{theorem}

\begin{theorem}[Uniform norm of a squared function]
\lean{FABL.uniformLpNorm_sq}
\label{FABL.uniformLpNorm_sq}
\leanok
\uses{FABL.uniformLpNorm, FABL.uniformLpNorm_abs_rpow}
\difficulty{2}
Let $\Omega$ be a finite set, let $q>0$, and let $f\colon\Omega\to\mathbb R$. The
uniform $L^{q/2}$ norm of the function $x\mapsto f(x)^2$ equals the square of the
uniform $L^q$ norm of $f$.
\end{theorem}

\begin{theorem}[Uniform norm of a constant function]
\lean{FABL.uniformLpNorm_const}
\label{FABL.uniformLpNorm_const}
\leanok
\uses{FABL.uniformLpNorm, FABL.uniformLpNorm_const_mul}
\difficulty{3}
For a finite nonempty set $\Omega$, a real number $q>0$, and a real constant $c$, the
uniform $L^q$ quantity of the constant function with value $c$ is $|c|$.
\end{theorem}

\begin{theorem}[Hypercontractivity from the last-coordinate parts]
\lean{FABL.twoQHypercontractivity_lastCoordinate_step}
\proofsource{Analysis-of-Boolean-Functions-by-Ryan-ODonnell}{
  pdf-ffbe9c7b3306-p273-b003,
  pdf-ffbe9c7b3306-p247-b006
}
\label{FABL.twoQHypercontractivity_lastCoordinate_step}
\leanok
\uses{FABL.Sign, FABL.lastEvenPart, FABL.lastFiberLpNorm, FABL.lastOddPart, FABL.noiseOperator, FABL.noiseOperator_snoc, FABL.oneBit_twoQ_sharp_of_two_le, FABL.signValue, FABL.uniformLpNorm, FABL.uniformLpNorm_add_le, FABL.uniformLpNorm_eq_lastFiberLpNorm, FABL.uniformLpNorm_le_of_abs_le, FABL.uniformLpNorm_nonneg, FABL.uniformLpNorm_sq, FABL.uniformLpNorm_sqrt_sq, FABL.uniformLpNorm_two_lastParts_sq}
\difficulty{6}
Let $n$ be a natural number, $q\ge2$, and $f$ a real-valued function on
$\{-1,1\}^{n+1}$. If the hypercontractive inequality with noise parameter $1/\sqrt{q-1}$
holds for both the average and the half-difference of the two restrictions of $f$
obtained by fixing the last coordinate to $1$ and $-1$, then it holds for $f$ itself.
\end{theorem}

\begin{theorem}[Two-norm hypercontractivity in every dimension]
\lean{FABL.twoQHypercontractivity_by_lastCoordinate}
\proofsource{Analysis-of-Boolean-Functions-by-Ryan-ODonnell}{
  pdf-ffbe9c7b3306-p273-b003,
  pdf-ffbe9c7b3306-p247-b006,
  pdf-ffbe9c7b3306-p259-b006
}
\label{FABL.twoQHypercontractivity_by_lastCoordinate}
\leanok
\uses{FABL.lastEvenPart, FABL.lastOddPart, FABL.noiseOperator, FABL.noiseOperator_apply_eq_pmfExpectation, FABL.pmfExpectation_const, FABL.twoQHypercontractivity_lastCoordinate_step, FABL.uniformLpNorm, FABL.uniformLpNorm_const}
\difficulty{4}
Let $q\ge2$. For every natural number $m$ and every real-valued function $f$ on
$\{-1,1\}^m$, applying the noise operator with parameter $1/\sqrt{q-1}$ to $f$ gives a
function whose uniform $L^q$ norm is at most the uniform $L^2$ norm of $f$.
\end{theorem}

\begin{theorem}[Noise operator contraction]
\lean{FABL.uniformLpNorm_noiseOperator_le}
\label{FABL.uniformLpNorm_noiseOperator_le}
\leanok
\uses{FABL.FiniteFourierBasis.productLpNorm_anisotropicProductNoiseOperator_le, FABL.anisotropicProductNoiseOperator_const, FABL.noiseOperator, FABL.productLpNorm_uniformSign_eq_uniformLpNorm, FABL.productNoiseOperator_uniformSign_eq_noiseOperator, FABL.signSiteFourierBasis, FABL.spectralProductNoiseOperator_eq_productNoiseOperator, FABL.uniformLpNorm}
\difficulty{3}
Let $n\ge 0$, $0\le\rho\le1$, and $p\ge1$. For every real-valued function $f$ on the
Boolean cube $\{−1,1\}^n$, the uniform $L^p$ norm of its noise-operator image at
parameter $\rho$ is at most the uniform $L^p$ norm of $f$.
\end{theorem}

\begin{theorem}[Monotonicity in the noise parameter]
\lean{FABL.uniformLpNorm_noiseOperator_mono}
\label{FABL.uniformLpNorm_noiseOperator_mono}
\leanok
\uses{FABL.noiseOperator, FABL.noiseOperator_comp, FABL.uniformLpNorm, FABL.uniformLpNorm_noiseOperator_le}
\difficulty{4}
Let $n\ge0$, let $0\le\rho\le\sigma\le1$, and let $p\ge1$. For every real-valued
function $f$ on $\{−1,1\}^n$, the uniform $L^p$ norm of its noise-operator image at
parameter $\rho$ is at most that of its image at parameter $\sigma$.
\end{theorem}

\begin{theorem}[Hypercontractivity sharp across two]
\lean{FABL.hypercontractivity_sharp_across_two}
\statementsource{Analysis-of-Boolean-Functions-by-Ryan-ODonnell}{
  pdf-ffbe9c7b3306-p258-b003,
  pdf-ffbe9c7b3306-p261-b003,
  pdf-ffbe9c7b3306-p273-b003
}
\label{FABL.hypercontractivity_sharp_across_two}
\leanok
\uses{FABL.noiseOperator, FABL.noiseOperator_comp, FABL.pTwoHypercontractivity, FABL.twoQHypercontractivity, FABL.uniformLpNorm}
\difficulty{4}
% dataset-pipeline: stub -- prose pending build_informal_statements.py
Statement pending regeneration.
\end{theorem}

\begin{theorem}[Hypercontractivity for the noise operator]
\lean{FABL.hypercontractivity_across_two}
\proofsource{Analysis-of-Boolean-Functions-by-Ryan-ODonnell}{
  pdf-ffbe9c7b3306-p273-b004,
  pdf-ffbe9c7b3306-p273-b003,
  pdf-ffbe9c7b3306-p258-b003
}
\label{FABL.hypercontractivity_across_two}
\leanok
\uses{FABL.hypercontractivity_sharp_across_two, FABL.noiseOperator, FABL.noiseOperator_comp, FABL.uniformLpNorm, FABL.uniformLpNorm_noiseOperator_le}
\difficulty{5}
Let $n\ge 0$, let $p,q,\rho\in\mathbb{R}$ satisfy $1\le p\le 2\le q$ and
$0\le\rho\le\sqrt{(p-1)/(q-1)}$, and let $f$ be a real-valued function on the sign cube
$\{\pm1\}^n$. Then the uniform $L^q$ norm of the noise operator with parameter $\rho$
applied to $f$ is at most the uniform $L^p$ norm of $f$.
\end{theorem}

\begin{theorem}[Square of a real power]
\lean{FABL.rpow_sq_eq_sq_rpow}
\label{FABL.rpow_sq_eq_sq_rpow}
\leanok
\difficulty{3}
For all real numbers $x,t$ with $x\ge0$, $(x^t)^2=(x^2)^t$, where real powers are used.
\end{theorem}

\begin{theorem}[Powers of the noise parameter]
\lean{FABL.noiseRpow_pow_sq}
\label{FABL.noiseRpow_pow_sq}
\leanok
\uses{FABL.rpow_sq_eq_sq_rpow}
\difficulty{3}
Let $n$ be a natural number, let $\rho,\lambda\in\mathbb{R}$ with $0\le\rho$, and let
$k$ be a natural number. Then
$\bigl((\rho^{1-\lambda})^k\bigr)^2=\bigl((\rho^k)^2\bigr)^{1-\lambda}$.
\end{theorem}

\begin{theorem}[Weighted geometric-sum inequality]
\lean{FABL.weighted_geometric_sum_le}
\label{FABL.weighted_geometric_sum_le}
\leanok
\difficulty{5}
Let $\iota$ be a type, let $s$ be a finite subset of $\iota$, and let $w_i,a_i\geq0$ for
every $i\in\iota$. For $0<\lambda<1$, one has $\sum_{i\in s}w_i a_i^{1-\lambda}\leq
(\sum_{i\in s}w_i)^\lambda(\sum_{i\in s}w_i a_i)^{1-\lambda}$, with powers interpreted
as real powers.
\end{theorem}

\begin{theorem}[Noise operator two interpolation]
\lean{FABL.noiseOperator_two_interpolation}
\statementsource{Analysis-of-Boolean-Functions-by-Ryan-ODonnell}{pdf-ffbe9c7b3306-p273-b004}
\label{FABL.noiseOperator_two_interpolation}
\leanok
\uses{FABL.fourierCoeff, FABL.fourierCoeff_noiseOperator, FABL.noiseOperator, FABL.noiseOperator_one, FABL.noiseRpow_pow_sq, FABL.parseval, FABL.rpow_sq_eq_sq_rpow, FABL.uniformLpNorm, FABL.uniformLpNorm_nonneg, FABL.uniformLpNorm_two_sq_eq_uniformInner, FABL.weighted_geometric_sum_le}
\difficulty{5}
% dataset-pipeline: stub -- prose pending build_informal_statements.py
Statement pending regeneration.
\end{theorem}

\begin{theorem}[Noise operator interpolation bound]
\lean{FABL.noiseOperator_q_interpolation}
\proofsource{Analysis-of-Boolean-Functions-by-Ryan-ODonnell}{pdf-ffbe9c7b3306-p273-b004}
\label{FABL.noiseOperator_q_interpolation}
\leanok
\uses{FABL.noiseOperator, FABL.noiseOperator_comp, FABL.noiseOperator_two_interpolation, FABL.twoQHypercontractivity, FABL.uniformLpNorm, FABL.uniformLpNorm_noiseOperator_mono}
\difficulty{5}
Let $n$ be a natural number, let $q,\rho,\lambda\in\mathbb{R}$ with $2\le q$,
$0\le\rho\le1$, and $0\le\lambda\le1$, and suppose $\rho^\lambda\le 1/\sqrt{q-1}$. For
every real-valued function $f$ on the $n$-dimensional sign cube, the uniform $q$-norm of
its image under the noise operator with parameter $\rho$ is at most the
$(1-\lambda)$-power of the uniform $2$-norm of that image times the $\lambda$-power of
the uniform $2$-norm of $f$.
\end{theorem}

\begin{theorem}[Reciprocal square-root inequality]
\lean{FABL.one_sub_le_inv_sqrt_one_add_two_mul}
\label{FABL.one_sub_le_inv_sqrt_one_add_two_mul}
\leanok
\difficulty{4}
For every $\epsilon\in[0,1]$, $1-\epsilon\le 1/\sqrt{1+2\epsilon}$.
\end{theorem}

\begin{theorem}[Bound for the noise operator on bounded functions]
\lean{FABL.abs_noiseOperator_le_one}
\label{FABL.abs_noiseOperator_le_one}
\leanok
\uses{FABL.abs_noiseOperator_le_noiseOperator_abs, FABL.noiseKernel, FABL.noiseOperator, FABL.noiseOperator_apply_eq_pmfExpectation, FABL.pmfExpectation, FABL.pmfExpectation_const, FABL.pmfExpectation_mono}
\difficulty{4}
Let $n\in\mathbb N$, let $\rho\in[0,1]$, and let $f\colon\{-1,1\}^n\to\mathbb R$ satisfy
$|f(y)|\leq 1$ for every $y\in\{-1,1\}^n$. Then for every $x\in\{-1,1\}^n$, $|T_\rho
f(x)|\leq 1$, where $T_\rho$ multiplies each Walsh coefficient indexed by a set $S$ by
$\rho^{|S|}$.
\end{theorem}

\begin{theorem}[Bounded noise q moment le]
\lean{FABL.bounded_noise_qMoment_le}
\statementsource{Analysis-of-Boolean-Functions-by-Ryan-ODonnell}{pdf-ffbe9c7b3306-p262-b003}
\label{FABL.bounded_noise_qMoment_le}
\leanok
\uses{FABL.abs_noiseOperator_le_one, FABL.noiseOperator, FABL.noiseOperator_comp, FABL.one_sub_le_inv_sqrt_one_add_two_mul, FABL.twoQHypercontractivity, FABL.uniformLpNorm, FABL.uniformLpNorm_noiseOperator_le, FABL.uniformLpNorm_nonneg, FABL.uniformLpNorm_rpow}
\difficulty{5}
% dataset-pipeline: stub -- prose pending build_informal_statements.py
Statement pending regeneration.
\end{theorem}

\begin{theorem}[Uniform expect prod le prod uniform lp norm]
\lean{FABL.uniform_expect_prod_le_prod_uniformLpNorm}
\statementsource{Analysis-of-Boolean-Functions-by-Ryan-ODonnell}{pdf-ffbe9c7b3306-p277-b002}
\label{FABL.uniform_expect_prod_le_prod_uniformLpNorm}
\leanok
\uses{FABL.uniformLpNorm, FABL.uniformLpNorm_nonneg, FABL.uniformLpNorm_rpow}
\difficulty{6}
% dataset-pipeline: stub -- prose pending build_informal_statements.py
Statement pending regeneration.
\end{theorem}

\begin{theorem}[Mixed even noise moment bound]
\lean{FABL.mixed_even_noise_moment_le}
\proofsource{Analysis-of-Boolean-Functions-by-Ryan-ODonnell}{pdf-ffbe9c7b3306-p277-b002}
\label{FABL.mixed_even_noise_moment_le}
\leanok
\uses{FABL.noiseOperator, FABL.twoQHypercontractivity, FABL.uniformLpNorm, FABL.uniformLpNorm_nonneg, FABL.uniformLpNorm_sq, FABL.uniformLpNorm_two_sq_eq_expect_sq, FABL.uniform_expect_prod_le_prod_uniformLpNorm}
\difficulty{4}
Let $r$ be a positive integer, let $n$ be a nonnegative integer, and let
$f_1,\ldots,f_r:\{-1,1\}^n\to\mathbb R$. Set $q=2r$ and $\rho=1/\sqrt{q-1}$. If $T_\rho$
denotes the noise operator that multiplies each Walsh coefficient of degree $k$ by
$\rho^k$, then, for uniform $x\in\{-1,1\}^n$, $$\mathbb E_x\prod_{i=1}^r (T_\rho
f_i(x))^2\leq \prod_{i=1}^r\mathbb E_x f_i(x)^2.$$
\end{theorem}
